\begin{abstract}
% Quantum computing promises exponential advantages on a number of classically
% intractable problems. However, quantum computers face large amounts of noise
% and must be accompanied by robust error correction to work properly. When
% running quantum computations, we cannot perform testing: unlike classical
% programs, quantum program outputs are costly to sample, may have non-classical
% output, and are produced on noisy hardware. Building end-to-end confidence in
% the correctness of quantum computations therefore requires verified
% error-correction guarantees at the foundation.
%
% We present a formalization of the theory of quantum error correction (QEC) in
% the Lean interactive theorem prover. Our library formalizes the linear algebra
% underlying quantum computing, the theory of the Pauli group, stabilizer codes,
% the binary symplectic representation, classical coding theory, and a number of
% well-studied QEC code families. A central contribution is an end-to-end
% framework for verifying code distance through an integrated SAT solver: code
% distance conditions are translated into Boolean satisfiability constraints via
% a verified reduction, exported to an external SAT solver, and the solver's
% result is reconstructed as a Lean-checked proof artifact. We demonstrate this
% pipeline on concrete code instances, including some from the practically
% relevant Bivariate Bicycle Code family, successfully verifying their distance
% properties within the proof assistant. We also demonstrate extensibility for
% larger code instances. Our work provides a reusable, machine-checked library
% for quantum error correction that fits into broader Lean-based quantum
% formalization efforts.
%
Quantum error-correcting codes (QECCs) sit between noisy quantum hardware and
reliable computation, so the code parameters used in practice must be
trustworthy. The single number that summarizes a code's strength is its
\emph{distance}, yet certifying a distance lower bound is NP-hard in general,
placing it beyond the reach of pen-and-paper proofs as well as direct
proof-assistant scripting. As a result, distance values in the literature
come either from non-scaling hand proofs, or from
unverified solvers that leave a trust gap exactly where the code is supposed to
provide a guarantee.
We present \textsc{Lean-QEC}, the first Lean 4 formalization of stabilizer-code
theory that delivers end-to-end, machine-checked distance certificates at
industrial code sizes. \textsc{Lean-QEC} formalizes the linear algebra of qubit
states, the Pauli group, stabilizer codes, the binary symplectic representation,
classical coding theory, and the CSS and Bivariate Bicycle families. To break the combinatorial barrier, \textsc{Lean-QEC} translates the distance
condition into a Boolean satisfiability formula through a verified reduction. The pipeline scales through a \textsc{BitVec}-flattened
encoding that replaces Lean's \textsc{Matrix} representation, and an
error-\emph{location} encoding that reduces the variable count from $n$ to
$k\lceil \log_2 n\rceil$. With these, we obtain automatically-generated Lean-checked distance proofs for a large range of industrially viable qLDPC codes within the Bivariate Bicycle and Generalized Bicycle families, including $\llbracket 90, 8, 10 \rrbracket$ and $\llbracket 70, 6, 9 \rrbracket$ BB codes, with the formulation scaling up to 144 qubits when performed outside the Lean kernel. The resulting library is reusable and is designed to
plug into broader Lean-based efforts toward end-to-end verification of
fault-tolerant quantum computation.
\end{abstract}
